\chapter{Análisis}
\label{chap:analisis}

\lettrine{E}{l} análisis del sistema parte de la identificación de los actores
que interaccionan con PadelCenter y de las necesidades que cada uno de ellos
tiene. A partir de estos actores se derivan los requisitos funcionales y no
funcionales, que constituyen la base del diseño y la implementación.

\section{Actores del sistema}
\label{sec:actores}

Se han identificado tres actores con responsabilidades diferenciadas, modelados
mediante el sistema de roles por centro (\texttt{CenterRole}):

\begin{description}
  \item[Usuario (\texttt{USER}):] Jugador registrado en el sistema. Puede
    consultar centros y pistas disponibles, crear y cancelar sus propias reservas,
    inscribirse en torneos como pareja y consultar sus partidos. Tiene visibilidad
    limitada: solo accede a sus propios datos.

  \item[Manager (\texttt{MANAGER}):] Gestor operativo de un centro concreto.
    Puede consultar todas las reservas del centro que gestiona, gestionar la
    disponibilidad de pistas y administrar los torneos del centro. Su ámbito de
    actuación está restringido al centro al que está asignado.

  \item[Administrador (\texttt{ADMIN}):] Administrador global del sistema.
    Tiene acceso completo: puede crear y eliminar centros, asignar roles a
    usuarios, gestionar cualquier entidad del sistema y acceder a todos los datos
    independientemente del centro.
\end{description}

El rol es contextual por centro: un mismo usuario puede ser \texttt{ADMIN} en
un centro y \texttt{USER} en otro. Esto se modela mediante la entidad
\texttt{CenterRole}, que asocia un usuario con un rol para un centro concreto.

\section{Requisitos funcionales}
\label{sec:requisitos_funcionales}

\subsection{Autenticación y usuarios}

\begin{itemize}
  \item \textbf{RF-01}: El sistema permite el registro de nuevos usuarios con
    nombre, apellidos, correo electrónico, contraseña y teléfono.
  \item \textbf{RF-02}: El sistema permite el inicio de sesión mediante
    credenciales (correo y contraseña), devolviendo un \textit{access token} JWT.
  \item \textbf{RF-03}: El \textit{access token} expira a los 15 minutos; el
    cliente puede renovarlo mediante un \textit{refresh token} almacenado en
    cookie \textit{httpOnly}.
  \item \textbf{RF-04}: Cada usuario tiene un rol global y opcionalmente uno
    o varios roles por centro (\texttt{CenterRole}).
  \item \textbf{RF-05}: Un administrador puede modificar los roles de cualquier
    usuario y cambiar su contraseña.
\end{itemize}

\subsection{Gestión de centros}

\begin{itemize}
  \item \textbf{RF-06}: Un administrador puede crear un centro deportivo con
    nombre, dirección, ciudad, teléfono y correo de contacto.
  \item \textbf{RF-07}: Un administrador puede eliminar un centro (soft delete),
    siempre que no tenga pistas activas asociadas; en caso contrario la
    operación se rechaza con un conflicto.
  \item \textbf{RF-08}: Cualquier usuario autenticado puede consultar el listado
    paginado de centros activos y el detalle de un centro.
\end{itemize}

\subsection{Gestión de pistas}

\begin{itemize}
  \item \textbf{RF-09}: Un administrador puede dar de alta pistas en sus centros,
    indicando nombre, tipo de superficie, precio por hora y disponibilidad.
  \item \textbf{RF-10}: Un administrador puede activar o desactivar una pista
    (\texttt{isAvailable}) sin eliminarla.
  \item \textbf{RF-11}: Cualquier usuario autenticado puede consultar la
    disponibilidad de una pista para una fecha concreta.
  \item \textbf{RF-12}: Un administrador puede eliminar una pista.
\end{itemize}

\subsection{Gestión de reservas}

\begin{itemize}
  \item \textbf{RF-13}: Un usuario puede crear una reserva para una pista en
    una franja horaria concreta.
  \item \textbf{RF-14}: El sistema rechaza la reserva si la pista no está
    disponible o si ya existe otra reserva activa en el mismo tramo horario
    (control de conflictos mediante índice de solapamiento en base de datos).
  \item \textbf{RF-15}: Una reserva pasa por los estados:
    \texttt{PENDING} → \texttt{CONFIRMED} → \texttt{COMPLETED}/\texttt{CANCELLED}.
  \item \textbf{RF-16}: Un usuario puede cancelar sus propias reservas; un
    administrador o manager puede cancelar cualquier reserva de su centro.
  \item \textbf{RF-17}: Un usuario puede consultar sus reservas actuales e
    históricas (\texttt{GET /api/v1/me/bookings}).
  \item \textbf{RF-18}: Un administrador o manager puede consultar todas las
    reservas de un centro (\texttt{GET /api/v1/centers/\{id\}/bookings}).
\end{itemize}

\subsection{Gestión de torneos}

\begin{itemize}
  \item \textbf{RF-19}: Un administrador puede crear un torneo en uno de sus
    centros con nombre, descripción, formato de competición y número máximo de
    parejas.
  \item \textbf{RF-20}: Los formatos de competición soportados son:
    \texttt{ROUND\_ROBIN} (todos contra todos), \texttt{ELIMINATION} (eliminación
    directa) y \texttt{GROUPS\_AND\_ELIMINATION} (fase de grupos más eliminación).
  \item \textbf{RF-21}: Un administrador puede abrir y cerrar el período de
    inscripción del torneo.
  \item \textbf{RF-22}: Los jugadores se inscriben en el torneo formando parejas.
    La pareja pasa por los estados \texttt{PENDING} → \texttt{CONFIRMED}/\texttt{REJECTED}.
  \item \textbf{RF-23}: Al cerrar la inscripción, el sistema genera automáticamente
    los emparejamientos del cuadro según el formato elegido.
  \item \textbf{RF-24}: Un administrador puede registrar el resultado de cada
    partido (\texttt{scoreA}, \texttt{scoreB}, ganador).
  \item \textbf{RF-25}: El sistema calcula y expone la clasificación en tiempo
    real para torneos en formato Round Robin.
  \item \textbf{RF-26}: Un usuario puede consultar sus torneos
    (\texttt{GET /api/v1/me/tournaments}) y sus partidos
    (\texttt{GET /api/v1/me/matches}).
\end{itemize}

\section{Ciclo de vida de un torneo}
\label{sec:ciclo_torneo}

El torneo sigue un ciclo de vida estricto con seis estados:

\begin{figure}[ht]
\centering
\resizebox{0.9\textwidth}{!}{%
\begin{tabular}{ccccc}
\texttt{DRAFT}
  & $\rightarrow$ &
\texttt{REGISTRATION\_OPEN}
  & $\rightarrow$ &
\texttt{REGISTRATION\_CLOSED} \\[0.4em]
  &   &   &   & $\downarrow$ \\[0.2em]
\texttt{CANCELLED}
  & $\leftarrow$ & \texttt{COMPLETED}
  & $\leftarrow$ & \texttt{IN\_PROGRESS} \\
\end{tabular}}

\vspace{0.4em}
\footnotesize\textit{\texttt{CANCELLED} es alcanzable desde cualquier estado anterior.}
\vspace{0.8em}

\renewcommand{\arraystretch}{1.25}
\begin{tabularx}{\textwidth}{lXX}
\hline
\textbf{Estado} & \textbf{Qué significa} & \textbf{Qué desencadena la transición de salida} \\
\hline
\texttt{DRAFT} &
  Borrador recién creado. No visible para inscripciones. Permite configurar el torneo antes de abrirlo. &
  El administrador invoca \textit{open registration}. \\
\texttt{REGISTRATION\_OPEN} &
  Las inscripciones están abiertas. Los jugadores pueden apuntarse como pareja. &
  El administrador cierra el período e invoca \textit{close \& generate bracket}. \\
\texttt{REGISTRATION\_CLOSED} &
  No se aceptan nuevas inscripciones. El sistema genera el cuadro de forma atómica al cerrar. &
  El administrador arranca el torneo (\textit{start}). \\
\texttt{IN\_PROGRESS} &
  Los partidos están en curso. Los managers registran resultados. La clasificación se actualiza en tiempo real. &
  Se completan todos los partidos previstos. \\
\texttt{COMPLETED} &
  Torneo finalizado. La clasificación final está disponible. Estado terminal. &
  — \\
\texttt{CANCELLED} &
  Torneo cancelado antes de completarse. No se aceptan más acciones sobre él. Estado terminal. &
  — \\
\hline
\end{tabularx}
\caption{Ciclo de vida de un torneo: estados, descripción y transiciones.}
\label{fig:ciclo_torneo}
\end{figure}

La transición \textit{close \& generate bracket} es atómica: cierra la
inscripción y genera todos los emparejamientos iniciales en una misma operación
({\ttfamily Close\-Registration\-And\-Generate\-Bracket\-Use\-Case}).

\section{Requisitos no funcionales}
\label{sec:requisitos_no_funcionales}

\begin{description}
  \item[RNF-01 — Seguridad:] Toda la API requiere autenticación mediante JWT
    Bearer token, salvo los endpoints de registro, login, refresh y logout. Las
    contraseñas se almacenan con BCrypt. El \textit{refresh token} se transmite
    exclusivamente mediante cookie \textit{httpOnly} y \texttt{SameSite=Strict}
    para mitigar ataques XSS y CSRF.

  \item[RNF-02 — Autorización granular:] La autorización se gestiona mediante
    evaluadores custom (\texttt{@bookingOwnerEvaluator},
    \texttt{@centerRoleEvaluator}, \texttt{@tournamentRoleEvaluator}) que se
    invocan desde \texttt{@PreAuthorize} en los casos de uso, garantizando que
    cada operación solo es accesible por el actor con el rol correcto en el
    contexto adecuado.

  \item[RNF-03 — Rendimiento:] El backend utiliza virtual threads de Java~21
    (\textit{Project Loom}), activados mediante
    \texttt{spring.threads.virtual.enabled=true}. Esto permite manejar un número
    mucho mayor de conexiones concurrentes sin saturar el pool de hilos del
    sistema operativo.

  \item[RNF-04 — Contrato API:] La especificación OpenAPI es la fuente de verdad.
    Cualquier cambio en la API debe reflejarse primero en el YAML y después
    propagarse al código generado tanto en backend como en frontend.

  \item[RNF-05 — Reproducibilidad:] El entorno de desarrollo es completamente
    reproducible mediante Docker Compose. Cualquier desarrollador puede levantar
    el stack completo con un único comando.

  \item[RNF-06 — Mantenibilidad:] La arquitectura hexagonal garantiza que el
    dominio de negocio no depende de ningún framework. Las reglas de arquitectura
    se verifican automáticamente en cada build mediante ArchUnit.

  \item[RNF-07 — Trazabilidad:] Todos los logs se emiten en formato JSON
    estructurado (Logstash encoder) con contexto MDC que incluye el identificador
    de usuario y el ID de correlación de la petición HTTP.
\end{description}

\section{Casos de uso}
\label{sec:casos_de_uso}

El sistema implementa 42 casos de uso organizados en cuatro categorías. La
tabla~\ref{tab:casos_de_uso} recoge los principales agrupados por dominio.

\begin{table}[ht]
\centering
\resizebox{\textwidth}{!}{%
\begin{tabular}{llp{7cm}}
\hline
\textbf{Dominio} & \textbf{Actor} & \textbf{Casos de uso} \\
\hline
\multirow{3}{*}{Auth / Usuario} & Anónimo & Registrarse, iniciar sesión, renovar token, cerrar sesión \\
 & Usuario & Ver perfil propio, actualizar contraseña \\
 & Admin & Listar usuarios, cambiar roles, eliminar usuario \\
\hline
\multirow{2}{*}{Centros} & Admin & Crear centro, eliminar centro \\
 & Todos & Listar centros, ver detalle de centro \\
\hline
\multirow{2}{*}{Pistas} & Admin/Manager & Crear pista, actualizar disponibilidad, eliminar pista \\
 & Todos & Listar pistas, ver disponibilidad por fecha \\
\hline
\multirow{3}{*}{Reservas} & Usuario & Crear reserva, cancelar reserva propia, ver mis reservas \\
 & Admin/Manager & Ver reservas del centro, cancelar cualquier reserva \\
 & Todos & Ver detalle de reserva \\
\hline
\multirow{4}{*}{Torneos} & Admin & Crear torneo, abrir inscripción, cerrar inscripción y generar cuadro, eliminar torneo \\
 & Admin/Manager & Confirmar pareja, reportar resultado de partido \\
 & Usuario & Inscribirse como pareja, ver mis torneos, ver mis partidos \\
 & Todos & Ver cuadro, ver partidos, ver clasificación (Round Robin) \\
\hline
\end{tabular}}
\caption{Casos de uso del sistema agrupados por dominio.}
\label{tab:casos_de_uso}
\end{table}

\subsection{Caso de uso detallado: Crear una reserva (CU-Booking-Create)}

\begin{description}
  \item[Actor principal:]\leavevmode Usuario autenticado con rol \texttt{USER}.
  \item[Precondición:]\leavevmode El usuario ha iniciado sesión, la pista existe y está
    activa (\texttt{isAvailable = true}).
  \item[Flujo principal:]\leavevmode
    \begin{enumerate}
      \item El usuario selecciona una pista e indica fecha, hora de inicio y
        hora de fin.
      \item El sistema verifica que la pista está disponible
        (\texttt{isAvailable}).
      \item El sistema consulta el índice de solapamiento en base de datos para
        detectar conflictos con reservas activas en el mismo tramo horario.
      \item El sistema crea la reserva con estado \texttt{PENDING}, calcula el
        precio total y la persiste.
      \item El sistema devuelve los datos de la reserva con HTTP 201.
    \end{enumerate}
  \item[Flujo alternativo — Conflicto de horario:]\leavevmode
    \begin{enumerate}
      \setcounter{enumi}{2}
      \item El sistema detecta solapamiento y devuelve HTTP 409 con detalle
        de error según RFC~9457.
    \end{enumerate}
  \item[Postcondición:]\leavevmode La reserva queda registrada con estado \texttt{PENDING}
    y la pista queda bloqueada para ese tramo horario.
\end{description}

\subsection{Caso de uso detallado: Cerrar inscripción y generar cuadro (CU-Tournament-CloseBracket)}

\begin{description}
  \item[Actor principal:]\leavevmode Administrador con rol \texttt{ADMIN} o \texttt{MANAGER}
    en el centro del torneo.
  \item[Precondición:]\leavevmode El torneo está en estado \texttt{REGISTRATION\_OPEN} y
    tiene al menos dos parejas confirmadas.
  \item[Flujo principal:]\leavevmode
    \begin{enumerate}
      \item El administrador invoca el endpoint de cierre de inscripción.
      \item El sistema cambia el estado del torneo a
        \texttt{REGISTRATION\_CLOSED}.
      \item El sistema genera automáticamente los emparejamientos iniciales
        según el formato del torneo (Round Robin, Eliminación o Grupos).
      \item Los partidos se crean con estado \texttt{SCHEDULED}.
      \item El sistema devuelve el cuadro generado con HTTP 200.
    \end{enumerate}
  \item[Postcondición:]\leavevmode El torneo tiene un cuadro completo y está listo para
    comenzar.
\end{description}
